\documentclass[11pt]{article}

\usepackage[utf8]{inputenc}
\usepackage[T1]{fontenc}
\usepackage{amsmath,amssymb,amsfonts}
\usepackage{graphicx}
\usepackage{booktabs}
\usepackage{hyperref}
\usepackage{natbib}
% \usepackage{lineno}
\usepackage{geometry}
\geometry{margin=1in}

\title{Hormonal-Social Adaptation Population Model:\\
A Computational Framework for Exploring Behavioral Stabilization Under Ecological Release}

\author{Jake Enholm\\
\textit{Data Scientist}\\
\small\url{https://github.com/jenholm/bio_models/tree/main/hsap}}

\date{2026}

\begin{document}

\maketitle

\begin{abstract}
Ecological conditions that reduce mortality pressure can alter the selective environment experienced by populations, yet most population regulation models focus on density-dependent resource competition and predation without incorporating endocrine-mediated behavioral feedbacks. We present HSAP (Hormonal-Social Adaptation Population model), an agent-based model testing whether coupled behavioral and endocrine mechanisms can influence population stability under reduced external threat. In HSAP, individuals regulate testosterone, cortisol, and estrogen in response to perceived predation risk and population density. These hormonal states drive behavioral outputs, including aggression, fertility, and mating drive, that feed back into population dynamics through a density-dependent behavioral sink mechanism with hysteresis. We evaluated 61 scenarios (4 core conditions, 3 sink variants, 54-factorial parameter sweep) across 1,970 validated simulation runs. Under reduced predation and disease pressure, populations stabilized at higher densities with reduced male testosterone and aggression (Cohen $d=3.29$). Under high space constraint with abundant resources, populations declined to 31\% of the resource-matched uncrowded condition, showing crowding pathology through fertility suppression. Under behavioral sink conditions, populations either recovered or produced partial collapse with 14\% extinction risk. A null model suite of 11 alternatives could not distinguish HSAP from logistic growth, predator-prey, or density-fertility models on population trajectories alone; HSAP became identifiable only when behavioral and endocrine observables were included. These results are consistent with the simulation assumptions: under the model's assumptions, endocrine-linked behavioral feedbacks can produce qualitatively distinct population dynamics under identical resource conditions, but population data alone are insufficient to identify the mechanistic pathways involved. The model does not establish that these mechanisms operate in natural populations; empirical validation requires paired endocrine and behavioral field measurements.
\end{abstract}

% \linenumbers

\section{Introduction}

\subsection{Population regulation through ecological pathways}

Population regulation is conventionally modeled through four classes of mechanism: resource competition, predation, disease, and density-dependent effects on fecundity and mortality \citep{sinclair1989,hixon2002population}. These frameworks, including logistic growth, Lotka-Volterra predator-prey dynamics, SIR disease models, and their extensions, share a common assumption: population size is determined by external conditions acting on individuals through resource limitation and mortality risk. While these mechanisms are well-established and empirically supported, they do not fully account for a recurring pattern in field observations: populations sometimes stabilize at densities well below resource-imposed carrying capacities, particularly following predator removal or during post-irruption phases \citep{krebs1995impact,krebs2001}.

\subsection{Behavioral regulation as an alternative pathway}

One class of explanations invokes behavioral plasticity. Individuals adjust aggression, reproductive effort, and social organization in response to perceived environmental conditions, creating feedback loops between individual physiology and population dynamics. This behavioral regulation pathway is distinct from resource-based regulation: it operates through individual decision rules rather than external mortality or fecundity limitation.

Endocrine hormones, particularly testosterone, cortisol, and estrogen, are plausible mediators of such adjustments. Testosterone influences aggression and reproductive investment across vertebrate taxa \citep{wingfield1990,sapolsky2005}. Cortisol mediates stress responses and can suppress reproductive function \citep{creel2013}. Estrogen regulates female reproductive effort and social behavior \citep{ellison2003}. The causal chain from environmental perception through hormonal state to behavioral output is well-documented in individual studies, but has rarely been incorporated into population-level models.

\subsection{Behavioral sink theory and crowding effects}

\citet{calhoun1962} identified ``behavioral sink'' phenomena in overcrowded populations, where density-dependent behavioral deterioration, including reproductive suppression, aggression, and social withdrawal, could trap populations in unsustainable states even after density was reduced. The behavioral sink concept suggests that populations can enter self-reinforcing states of reproductive failure that persist beyond the ecological conditions that triggered them.

Subsequent work has extended this framework. \citet{christian1964endocrines} proposed endocrine-mediated mechanisms linking population density to reproductive suppression through stress hormone pathways. These studies converge on a common hypothesis: that behavioral and endocrine feedbacks can generate population-level regulatory effects independent of resource competition.

\subsection{Existing models and the integration gap}

Several modeling frameworks have explored aspects of endocrine-mediated population regulation. Endocrine-based models of stress and dominance have been applied to primate social dynamics \citep{sapolsky2005}, but typically at the group level rather than as population-level regulators. Agent-based models of behavioral ecology have incorporated hormonal state variables \citep{higham2013endocrinology}, but usually without formal null model comparisons to test whether endocrine mechanisms add explanatory power beyond simpler alternatives.

The gap in the literature is a model that integrates: (1) individual-level endocrine feedback from environment to behavior; (2) sex-specific behavioral dynamics (male aggression, female competition); (3) a density-driven behavioral sink with recovery; and (4) rigorous null model comparison to test whether population trajectories alone can identify the endocrine mechanism.

\subsection{HSAP contribution}

We present HSAP (Hormonal-Social Adaptation Population model) as a falsifiable agent-based model that addresses this gap. HSAP proposes a specific causal chain: low external threat $\rightarrow$ male testosterone downshift $\rightarrow$ reduced male aggression, elevated female aggression $\rightarrow$ reduced fertility $\rightarrow$ population stabilization. The model is designed to be tested against its own predictions, not merely to show that endocrine feedbacks can generate such dynamics.

Our contributions are: (1) a computational framework linking individual endocrine state to population dynamics through behavioral feedbacks; (2) evidence that identical resource conditions produce qualitatively different population outcomes depending on behavioral and endocrine parameters; (3) an honest null model comparison showing that population trajectories alone are insufficient for mechanism identification; and (4) identification of the observable classes required for empirical testing.

\textbf{Scope and boundaries.} HSAP is a mammalian population model. It does not model human sexual orientation; non-reproductive sexual or social behavior is treated only as species-specific ethological behavior requiring explicit supporting data. The model is not empirically validated; empirical validation requires paired endocrine and behavioral field measurements.

\begin{figure}[htbp]
\centering
\includegraphics[width=0.95\linewidth]{figures/figure_1_causal_chain.pdf}
\caption{Conceptual structure of the HSAP model. Ecological conditions influence endocrine-linked behavioral state variables, which feed back into fertility, aggression, mortality, and population dynamics.}
\label{fig:causal_chain}
\end{figure}



\section{Methods}

\subsection{Agent Architecture}

Each agent in HSAP is characterized by sex (assigned 50/50 at birth), age (in model time steps; 10 steps = 1 model-year), and a suite of physiological and behavioral state variables. Agents are initialized with age drawn uniformly from $[5, 40)$ and energy, health, and rank set to 1.0, 1.0, and 0.5 respectively.

\paragraph{Hormonal state.}
Each agent carries three hormonal axes: testosterone, estrogen, and cortisol. Male testosterone is initialized near 1.0 (\texttt{base\_male\_testosterone}) with Gaussian noise (SD=0.15); female testosterone near 0.3 (SD=0.05). Estrogen is initialized near 1.0 for females and 0.3 for males. Cortisol is initialized near 0.5 for both sexes (SD=0.1). All hormonal values are clamped to $[0, 1]$ at each time step.

\paragraph{Behavioral state.}
Derived from hormones and environment at each step: \texttt{aggression\_tendency}, \texttt{maternal\_defense\_tendency} (females), \texttt{mating\_drive}, \texttt{fertility}, and \texttt{offspring\_survival\_probability}. These are recomputed every time step; agents do not carry fixed behavioral types.

\paragraph{Social state.}
Each agent has a dominance rank (clamped $[0, 1]$), territory quality, \texttt{cooperation\_bonus} (accumulated from cooperative actions, capped at 0.3), and a pregnancy flag with gestation countdown (\texttt{gestation\_time} = 5 steps).

\subsection{Environment}

The environment provides four static parameters and several dynamic state variables.

\paragraph{Static parameters} (set per scenario):
\begin{itemize}
\item \texttt{resource\_abundance} (0 to 1): baseline resource availability
\item \texttt{predator\_pressure} (0 to 1): external threat level
\item \texttt{disease\_pressure} (0 to 1): disease prevalence
\item \texttt{space\_constraint} (0 to 1): controls how quickly territory availability declines with density
\item \texttt{carrying\_capacity} (integer): $K$ for density computation
\end{itemize}

\paragraph{Dynamic state} (updated each step):
\begin{itemize}
\item \texttt{resource\_abundance}: regenerates toward maximum, consumed by population; modulated by seasonality
\item \texttt{territory\_availability}: $\max(0, 1 - \rho \cdot S)$, where $\rho = N / K$
\item \texttt{predation\_risk}: $P \cdot (1 + 0.5 \cdot \rho)$
\item \texttt{disease\_risk}: $D \cdot (1 + 0.5 \cdot \rho)$
\end{itemize}

\subsection{Endocrine Update}

Hormones are recomputed for every agent at each time step, creating the causal feedback loop that defines HSAP.

\paragraph{Male testosterone} is driven by five factors:

\begin{equation}
T_{m,t+1} = \min\left(T_{\max},\; T_{\text{base}} + 0.3 \cdot \text{rank} + 0.4 \cdot M_d - 0.3 \cdot \rho - 0.4 \cdot (1 - P) - 0.2 \cdot (1 - \text{health})\right)
\label{eq:male_T}
\end{equation}

The term $-0.4 \cdot (1 - P)$ is the critical endocrine feedback: when predator pressure is high, the term is near zero (testosterone is not reduced); when predator pressure is low, the term subtracts up to 0.4 (testosterone is downshifted). This is the mechanism through which threat reduction propagates to behavioral change.

\paragraph{Female testosterone} responds to resource scarcity and density:

\begin{equation}
T_{f,t+1} = \min\left(T_{\max},\; 0.3 + 0.15 \cdot \text{rank} + 0.4 \cdot (1 - R) + 0.3 \cdot \rho\right)
\label{eq:female_T}
\end{equation}

\paragraph{Cortisol} (both sexes) integrates density, rank, health, predation, and resource scarcity:

\begin{equation}
C = 0.5 + 0.1 \cdot \rho - 0.3 \cdot \text{rank} + 0.5 \cdot (1 - \text{health}) + 0.2 \cdot P + 0.3 \cdot (1 - R)
\label{eq:cortisol}
\end{equation}

\paragraph{Estrogen} (females) is suppressed by cortisol:

\begin{equation}
E = 1.0 - 0.3 \cdot C
\label{eq:estrogen}
\end{equation}

\subsection{Behavioral Output}

From hormones, agents derive:

\begin{itemize}
\item \textbf{Aggression}: base $= T \cdot 0.5 + C \cdot 0.2$, with sex-specific additions (females: offspring defense + resource competition; males: dominance bonus)
\item \textbf{Fertility}: base $= 0.5$, modified by hormones, age, density-dependent reproductive restraint (activates above density 0.7), and territory availability
\item \textbf{Mating drive}: males $= T \cdot 0.5 + 0.3$; females $= E \cdot 0.4 + 0.2$
\item \textbf{Offspring survival probability}: affected by testosterone, cortisol, density, and resource scarcity
\end{itemize}

Agents select actions stochastically from: forage, dominate, disperse, withdraw, cooperate, guard\_offspring. Action probabilities depend on current behavioral state and environmental conditions.

\subsection{Reproduction}

Mating occurs when a randomly selected male (probability proportional to mating\_drive) encounters a fertile female above a fertility gate (default 0.3). Both must have energy above a threshold (0.4) and be in reproductive age (10--60 model-years). On successful mating, the female becomes pregnant with gestation time of 5 steps.

At birth, litter size is drawn from a normal distribution (mean=4.0, SD=1.0, floored to 1). Each pup survives with probability determined by the mother's \texttt{offspring\_survival\_probability}. The mother pays an energy cost of 0.3. Newborns face an additional mortality check.

\paragraph{Infanticide} triggers when density exceeds 0.8, with probability proportional to agent aggression and density above threshold. \paragraph{Offspring neglect} triggers when density exceeds 0.7 or cortisol exceeds 0.7, removing one offspring.

\subsection{Mortality}

Death probability per time step is computed as:

\begin{equation}
P_{\text{death}} = \min\left(1,\; (P_{\text{base}} + P_{\text{pred}} + P_{\text{disease}} + P_{\text{starvation}} + P_{\text{injury}} + P_{\text{old}} + P_{\text{crowd}}) \times \tau \times M\right)
\label{eq:mortality}
\end{equation}

where $\tau = 0.1$ (age increment per step) and $M$ is the mortality multiplier (1.0 normally, 0.4 during post-sink recovery).

\begin{table}[h]
\centering
\caption{Mortality components}
\label{tab:mortality}
\begin{tabular}{ll}
\toprule
Component & Formula \\
\midrule
Base & 0.02 (constant) \\
Predation & $0.05 \cdot P_{\text{risk}} \cdot (1 - 0.02 \cdot (1-R)) \cdot \text{modifiers}$ \\
Disease & $0.03 \cdot D_{\text{risk}} \cdot (1 + 0.5 \cdot (1 - \text{health})) \cdot (1 + 0.5 \cdot \rho)$ \\
Starvation & If energy $< 0.2$: $0.3 \cdot (1-R) \cdot (1 - \text{energy}/0.2)$ \\
Injury & 0.1 if injured \\
Old age & $0.01 \cdot (\text{age} - 50)$ if age $> 50$ \\
Crowding & $0.05 \cdot \max(0, \rho - 0.8) + \eta \cdot M_{\text{sink}}$ \\
\bottomrule
\end{tabular}
\end{table}

\subsection{Behavioral Sink and Recovery}

The behavioral sink is a density-driven phase transition that suppresses reproductive output even after density falls below carrying capacity.

\paragraph{Activation.}
When density exceeds \texttt{behavioral\_sink\_on\_threshold} (default 0.75), the sink engages. A \texttt{sink\_factor} is computed as a continuous measure of sink severity:

\begin{equation}
\eta = \min\left(1,\; \frac{\rho - 0.75}{1 - 0.75}\right)
\label{eq:sink_factor}
\end{equation}

\paragraph{Effects while active.}
Fertility is multiplicatively suppressed: $F \times (1 - \eta \cdot 0.85)$. Mating drive is suppressed: $M_d \times (1 - \eta \cdot 0.4)$. Withdrawal behavior and offspring neglect probabilities are increased.

\paragraph{Deactivation (hysteresis).}
The sink disengages only when density falls below \texttt{behavioral\_sink\_off\_threshold} (default 0.50), creating a hysteresis band (on at 0.75, off at 0.50). A minimum duration of 30 steps must elapse before deactivation is permitted. An optional auto-recovery timer forces disengagement after a fixed number of steps.

\paragraph{Post-sink recovery.}
After sink disengagement, a recovery phase lasting 100 steps begins. During recovery: fertility and mating drive are multiplicatively boosted by 1.3; mortality is reduced to 40\% of normal. If population drops below 30 during recovery, refugees (up to 10 per step) are injected to supplement the depleted population.

\subsection{Metrics}

Per time step, HSAP records: population size, density, sex ratio, resource abundance, predator pressure, disease pressure, territory availability, age structure (juvenile/adult/senior counts), mean hormones (testosterone, estrogen, cortisol, male\_T, female\_T), behavioral outputs (male/female aggression, female defense), reproductive metrics (fertility, births, deaths, matings, pregnancies, infanticide, neglect, refugees), sink state (active, post-recovery, sink\_factor), external threat index, and the composite HSAP index.

The HSAP index is a weighted composite of five signals: low external threat (25\%), low male aggression (20\%), high female aggression (20\%), low fertility (20\%), and population stability (15\%), with a viability dampener for populations below 50.

\paragraph{Summary metrics} (computed over the full run): final population, peak population, minimum population, crash ratio (final/peak), time to stability, mean male/female testosterone, mean male/female aggression, mean fertility, mean cortisol, extinction status, and seed number.

\section{Results}

%% TODO: Verify all numbers against frozen data
%% TODO: Check all CI calculations

\subsection{Core Scenario Separation}

All four core scenarios produced distinct population outcomes (Table~\ref{tab:core}, Figure~\ref{fig:trajectories}).

\textbf{A\_normal\_baseline} stabilized at a mean population of 94 agents (95\% CI=$[89, 100]$) with zero extinctions across 50 seeds. Male testosterone averaged 0.994 (CI=$[0.993, 0.995]$).

\textbf{B\_hsap\_abundance} (reduced threat) stabilized at 141 agents (CI=$[136, 147]$), a 50\% increase over baseline. Male testosterone shifted to 0.986 (CI=$[0.985, 0.987]$), directionally consistent with relaxed competitive endocrine signaling. The 95\% confidence intervals for male testosterone do not overlap between A and B (Cohen $d=3.29$). Male aggression similarly decreased from 0.895 (CI=$[0.892, 0.897]$) to 0.847 (CI=$[0.846, 0.849]$). The effect was statistically robust but modest in absolute magnitude (0.8 percentage points on a $[0,1]$ scale), indicating that the model's population-level separation is not driven by a large endocrine swing alone.

%% TODO: Note that Cohen d is inflated by simulation precision — consider reporting effect size with biological context

\textbf{C\_crowded\_abundance} (space constraint) stabilized at 44 agents (CI=$[41, 46]$), 47\% below baseline despite identical resource availability. Fertility was reduced to 0.436 (CI=$[0.433, 0.440]$) compared to 0.568 in B, indicating crowding pathology. Male testosterone was 0.982 (CI=$[0.980, 0.984]$), slightly lower than B, consistent with density-mediated endocrine suppression.

\textbf{D\_high\_predation\_survival} (elevated threat) stabilized at 77 agents (CI=$[72, 81]$), 18\% below baseline. Male testosterone (0.997, CI=$[0.997, 0.998]$) and aggression (0.933, CI=$[0.932, 0.935]$) were elevated above baseline, consistent with predator-mediated selection pressure on endocrine state.

\subsection{Crowding Pathology}

The comparison between B\_hsap\_abundance (pop=141, fertility=0.568) and C\_crowded\_abundance (pop=44, fertility=0.436) shows that space constraint alone, with identical resources, can suppress population to 31\% of the resource-matched uncrowded condition. The crowding pathology manifests through reduced fertility, not increased mortality, consistent with density-dependent reproductive restraint rather than direct population control.

\subsection{Sink, Recovery, and Collapse}

The behavioral sink produced qualitatively distinct dynamics depending on scenario conditions (Table~\ref{tab:sink}, Figure~\ref{fig:sink}).

\textbf{C\_crowded\_stable} (control) produced outcomes identical to C\_crowded\_abundance: pop=44, male\_T=0.982. This is consistent with the behavioral sink requiring its density-triggered mechanism; removing the trigger eliminates sink dynamics.

\textbf{E\_behavioral\_sink\_recovery} produced the characteristic sink-recovery pattern. The sink engaged in 62\% of timesteps and disengaged in 36\%. Population dipped from an initial transient to a minimum of $\sim$75 agents, then recovered to $\sim$98 agents. The behavioral sink is reversible under these conditions: the hysteresis mechanism allowed density to fall below the deactivation threshold, triggering recovery with fertility and mating drive boosts and mortality relief.

\textbf{F\_behavioral\_sink\_partial\_collapse} produced severe partial collapse. Population was reduced to $\sim$4 agents (CI=$[3, 5]$). Sink engagement rate was 55\%, but recovery rate was 0\% --- the behavioral sink never disengaged. Extinction occurred in 14\% of seeds (95\% CI=$[6\%, 24\%]$). This is partial collapse with nonzero extinction risk. The 14\% extinction rate reflects stochastic variation in mortality and reproduction under conditions of sustained reproductive suppression.

\subsection{Null Model Non-Identifiability}

The null model comparison reveals that population trajectories alone are partly non-identifiable across the null model suite (Table~\ref{tab:null}, Figure~\ref{fig:null}).

Of 11 null models tested, 7 (logistic growth, predator-prey, density-fertility, dominance hierarchy, random hormone, density with recovery, random behavior) could not distinguish A from B using population trajectories alone (A/B MSE ratio $< 1.5$). Two models (disease pressure, endocrine-no-behavior) showed partial separation but with MSEs 25,000--86,000 compared to HSAP ablation MSEs of 70--2,458. One model (resource-only) distinguished A from B but with poor overall fit (MSE=23,850 for A).

%% TODO: Justify the 1.5x MSE threshold — is this arbitrary?

This non-identifiability is one of the paper's main results. It implies that a field study measuring population size alone cannot reject HSAP in favor of logistic growth, predator-prey dynamics, or density-dependent fertility. HSAP becomes testable only when population outcomes are paired with behavioral and endocrine observables --- specifically, male testosterone and aggression measurements that null models cannot reproduce.

\subsection{Ablation Results}

The HSAP ablation analysis quantifies the contribution of each model component (Table~\ref{tab:ablation}, Figure~\ref{fig:ablation}).

Removing the male testosterone downshift channel increased B's population MSE from 73 (full model) to 565 --- a 7.8-fold increase. Removing the female aggression channel increased it to 397 (5.4-fold). These two components are the minimal endocrine mechanisms required for HSAP's distinguishing power: they produce population separation between A and B while maintaining low overall MSE.

Removing cortisol increased MSE to 772--898 across scenarios, indicating a broader but less scenario-specific effect. Removing endocrine responsiveness (the coupling between hormones and behavior) produced larger MSE increases on B (2,458) than on A (913), indicating that the endocrine-behavioral coupling is more consequential under reduced threat.

Removing sink recovery produced MSE=0 across all scenarios, because without the sink mechanism, the behavioral sink never engages and the recovery phase never begins. This is consistent with the internal consistency of the model.

\subsection{Factorial Sensitivity}

The 54-factorial analysis identifies space constraint as the dominant population regulator (Table~\ref{tab:factorial}, Figure~\ref{fig:phase}). Under low space constraint (0.40), populations ranged from 99--143 agents regardless of predation, disease, or resource levels. Under high space constraint (0.80), populations ranged from 4--13 agents under high predation and disease, showing that crowding pathology can override resource abundance.

The factorial analysis also reveals interaction effects: high predator pressure combined with high disease pressure produced population collapse only under high space constraint, suggesting that multiple stressors are necessary for collapse when space is abundant.

\section{Discussion}

\subsection{What HSAP Supports}

Under model assumptions, HSAP shows that endocrine-linked behavioral feedbacks can generate qualitatively distinct population outcomes under identical resource conditions. The same resource abundance (0.60) produced populations of 141 (B\_hsap\_abundance), 44 (C\_crowded\_abundance), and 77 (D\_high\_predation\_survival), depending on threat level and space constraint. These differences are simulation-consistent with the hypothesis that endocrine state mediates population regulation, not just resource availability.

The behavioral sink mechanism produces two distinct outcomes depending on conditions: recovery (E) or partial collapse (F). The difference is not in the mechanism itself but in the environmental context --- F operates under higher predation and disease pressure, preventing the density reduction needed to disengage the sink. This is simulation-consistent with the conservation concept that multiple stressors can trap populations in unsustainable behavioral states.

\subsection{What HSAP Does Not Prove}

HSAP is a simulation model, not an empirical study. It does not prove that endocrine feedbacks drive population dynamics in any specific species. It does not validate the specific parameter values used (which are illustrative, not empirically measured). It does not demonstrate that the behavioral sink mechanism exists in nature.

The model shows that IF endocrine-linked behavioral feedbacks exist AND operate as modeled, THEN population-level consequences follow. The ``if'' is the empirical question that field studies must address.

\subsection{Why Endocrine and Behavioral Measurements Are Necessary}

The null model comparison is the paper's most important result for empirical science. Seven of 11 null models could not distinguish HSAP from simpler alternatives on population trajectories alone. Population data alone are insufficient to identify the mechanistic pathways driving stabilization or collapse.

However, HSAP generates testable predictions in the joint space of population size, endocrine state, and behavior. A field study measuring population size AND male testosterone AND aggression can reject null models while retaining HSAP. The behavioral and endocrine observables that distinguish HSAP --- particularly the male testosterone downshift under reduced threat --- produce separable signatures in the observable variable space that no single-channel null model reproduces.

This is a call for richer data in conservation biology, not a dismissal of population monitoring. Population data remain essential, but they are insufficient for mechanism identification.

\subsection{Empirical Predictions}

HSAP generates three testable predictions for post-predator-removal populations:

\begin{enumerate}
\item \textbf{Endocrine signature.} Male testosterone should decrease following predator removal, with the magnitude proportional to the reduction in perceived threat. This is the most direct test of the endocrine feedback loop.

\item \textbf{Crowding pathology.} Populations in space-constrained environments should show higher cortisol and lower fertility than resource-matched populations in unconstrained environments, even when resource availability is identical.

\item \textbf{Behavioral sink persistence.} Populations under high combined stressors (predation + disease + crowding) should show persistent reproductive suppression despite below-capacity density, consistent with the behavioral sink mechanism.
\end{enumerate}

These predictions are falsifiable and could be tested in post-predator-removal monitoring programs that include endocrine sampling.

\subsection{External Endocrine and Behavioral Proxy Evidence}

The external proxy evidence should be interpreted as empirical motivation rather than validation. Published secular endocrine studies report declines in male testosterone over multi-decade timescales \citep{travison2007,lokeshwar2020}. Sex-disaggregated arrest statistics show increased female share of assault-related arrests relative to male trends. These trends are directionally compatible with the paper's methodological claim that HSAP-like mechanisms require paired endocrine and behavioral observables for empirical testing: the endocrine trend is consistent with relaxed competitive signaling under reduced threat, and the behavioral trend is consistent with altered behavioral output under crowding and sink conditions.

However, the available series are ecological and confounded. Endocrine studies are cross-sectional and do not track individual-level change. Arrest data reflect behavior filtered through reporting, enforcement, and charging practices. The two data streams are not linked at the individual level. Their main value is to show how HSAP-like mechanisms could be tested: by pairing endocrine measurements with behavioral proxies rather than relying on population outcomes alone. Figure~\ref{fig:observable} illustrates the observable requirements: population data alone are non-identifiable, but endocrine and behavioral observables together define the variable space needed for empirical testing. Contemporary endocrine and behavioral proxy trends identify the class of observables needed to test HSAP-like mechanisms outside simulation, but these ecological data are exploratory and non-causal.

%% TODO: Need citation for Travison 2007 secular testosterone decline
%% TODO: Need citation for Lokeshwar 2020
%% TODO: Need to ensure human implications are explicitly avoided per README boundary

\section{Conclusion}
\label{sec:conclusion}

HSAP is a falsifiable agent-based model showing that endocrine-linked behavioral feedbacks can generate population stabilization, crowding pathology, recovery, or partial collapse under reduced external threat. The model's most important finding is negative: population trajectories alone are insufficient to identify the mechanistic pathways driving these dynamics. Seven of 11 null models produced indistinguishable population outcomes. HSAP becomes testable only when population outcomes are paired with behavioral and endocrine observables.

This result has practical implications for conservation biology. Post-predator-removal monitoring programs that measure only population size may miss endocrine and behavioral signatures of regulatory change. HSAP argues for integrating endocrine sampling into population monitoring, particularly in systems where predator removal has produced unexpected population dynamics.

The model is not empirically validated. It shows that endocrine-linked behavioral feedbacks are sufficient to produce qualitatively distinct population dynamics under identical resource conditions. Whether these feedbacks are necessary, or even present, in real populations remains an empirical question.


\bibliographystyle{plainnat}
\bibliography{references}

\end{document}
